% Pista alg-25j-q2 · Àlgebra
% Procedència: PAU Catalunya 2025 · Convocatòria de juny · Sèrie 1 · Exercici 2
\documentclass[11pt,a4paper]{article}
\usepackage{pista}
\pistainfo{Catalunya 2025}{Convocatòria de juny}{Sèrie 1}

\begin{document}

\section*{\textcolor{blauPista}{Exercici 2}}

\noindent Considereu el sistema d'equacions lineals següent:
\[
\left.\begin{array}{r}
y - z = p + 3 \\
p^{2}x - z = 5 \\
x - y = 3
\end{array}\right\}
\]

\subsection*{\textit{a)} \normalfont Discutiu el sistema per als diferents valors del paràmetre $p$. \hspace*{\fill}\textnormal{[1,25~punts]}}

\pista{Abans de calcular res, escriu $A$ i $A^{'}$ amb les incògnites en l'ordre $x$, $y$, $z$. On falta una incògnita, hi va un $0$. Per exemple, la primera equació, $y - z = p + 3$, dona la primera fila d'$A^{'}$:
\[
\left(\begin{array}{ccc|c} 0 & 1 & -1 & p + 3 \end{array}\right)
\]
A continuació, estudia el sistema d'equacions. D'una banda, has d'estudiar a part cada valor de $p$ que anul·la el determinant d'$A$. D'altra banda, la resta de valors que no anul·len el determinant fan que el sistema sigui compatible determinat, és a dir, que té solució única.}

\subsection*{\textit{b)} \normalfont Resoleu el sistema per al cas $p = -1$. \hspace*{\fill}\textnormal{[0,5~punts]}}

\pista{Substitueix $p = -1$ a les tres equacions.

\medskip
Mira què t'ha sortit a l'apartat \textit{a)} per a $p = -1$. Si el sistema és compatible indeterminat, en tens prou amb dues equacions: fes servir la 1a i la 3a. A continuació, escriu simplement \mbox{$y = t$} (és un paràmetre lliure), i finalment aïlla $z$ de la 1a equació i $x$ de la 3a.}

\subsection*{\textit{c)} \normalfont Per al cas $p = -1$, hi ha alguna solució que compleixi, a més, $xy = 10$? En cas afirmatiu, indiqueu quantes n'hi ha i trobeu-les totes. \hspace*{\fill}\textnormal{[0,75~punts]}}

\pista{En primer lloc, substitueix les expressions d'$x$ i d'$y$ que has trobat a l'apartat \textit{b)} a la condició $x \cdot y = 10$. Obtindràs una equació de segon grau amb la variable $t$.

\medskip
Un cop hagis resolt aquesta equació, per a cada valor obtingut de $t$, calcula la solució completa $(x, y, z)$, i digues quantes n'hi ha.}

\vspace{0.6em}

\end{document}
